Une usine hydroélectrique fonctionne grâce à une chute d'eau. Le débit de l'eau
est de $\qty{8.0}{\cubic\m\per\s}$ et la hauteur de chute est
de $\qty{20}{\m}$.
\begin{enumerate}
    \item Calculer la puissance de cette chute d'eau sachant qu'elle est égale
          au travail fourni en $\qty{1.0}{\s}$ par le poids de l'eau subissant
          la chute.
    \item Quel est le travail fourni en 1 jour ($\qty{24}{\h}$)~?
\end{enumerate}

\begin{donnees}
    \begin{itemize}
        \item $g=\qty{10}{\N\per\kg}$.
        \item masse volumique de l'eau~: $\rho=\qty{1000}{\kg\per\cubic\m}$.
    \end{itemize}
\end{donnees}

